\chapter{방정식에서 현대대수학까지}

\begin{learninggoal}
이 장을 마치면 다음을 설명할 수 있다.
\begin{itemize}
  \item 방정식의 해법 연구가 왜 대수적 구조의 연구로 전환되었는가?
  \item 근의 공식과 ``근호로 풀 수 있음''은 정확히 무엇을 뜻하는가?
  \item 근들의 치환이 왜 군이라는 구조를 이루는가?
  \item 군론, 환론, 체론이 갈루아 이론 안에서 어떻게 연결되는가?
\end{itemize}
\end{learninggoal}

\section{계산법이 아니라 가능성을 묻기}

일차방정식과 이차방정식에는 일반적인 해법이 있다. 삼차와 사차방정식에도 계수의 사칙연산과 근호를 이용한 공식이 존재한다. 자연스럽게 다음 질문이 생긴다.

\begin{quote}
모든 차수의 다항방정식에 이와 같은 공식이 존재하는가?
\end{quote}

이 질문의 핵심은 특정 방정식 하나를 푸는 것이 아니다. 주어진 차수의 모든 방정식에 적용되는 \emph{일반적인 절차}가 존재하는지를 묻는 것이다. 오차방정식에 대한 연구는 결국 ``공식을 찾아라''에서 ``그런 공식이 존재할 수 있는가''로 바뀌었다.

이 전환은 현대대수학의 전형적인 사고방식이다. 계산을 더 길게 수행하는 대신, 계산이 가능하게 만드는 구조를 먼저 찾는다.

\section{근호로 풀 수 있다는 말}

계수가 유리수인 다항식
\[
 f(x)=a_nx^n+\cdots+a_1x+a_0
\]
을 생각하자. 이 방정식이 근호로 풀린다는 것은, 유리수에서 시작해 사칙연산과 유한 번의 $m$제곱근 추출을 반복하여 모든 근을 표현할 수 있다는 뜻이다.

예를 들어
\[
 x^2-2=0
\]
의 근은 $\pm\sqrt{2}$이므로 근호로 풀린다. 반면 일반적인 오차방정식에는 모든 경우에 적용되는 근호 공식이 존재하지 않는다. 중요한 점은 ``어떤 오차방정식도 풀 수 없다''가 아니라, 오차방정식 전체에 통용되는 하나의 근호 공식이 없다는 사실이다.

\begin{pitfall}
``일반 오차방정식은 근호로 풀 수 없다''와 ``모든 오차방정식은 근호로 풀 수 없다''를 혼동하면 안 된다. $x^5-1=0$처럼 근호로 풀리는 오차방정식은 많다.
\end{pitfall}

\section{근들의 대칭}

다항식의 계수만 알고 있을 때, 근에는 이름표가 붙어 있지 않다. 근을 $\alpha_1,\ldots,\alpha_n$이라 적더라도 이 번호는 우리가 임의로 정한 것이다. 따라서 근들을 서로 바꾸어도 계수로 표현되는 관계가 보존될 수 있다.

예를 들어 이차방정식
\[
 x^2-sx+p=0
\]
의 두 근을 $\alpha,\beta$라 하면
\[
 \alpha+\beta=s,\qquad \alpha\beta=p
\]
이다. $\alpha$와 $\beta$를 맞바꾸어도 두 식은 변하지 않는다. 이 단순한 교환은 근들 사이의 대칭이다.

근들의 치환을 연달아 적용할 수 있고, 아무것도 바꾸지 않는 치환이 있으며, 각 치환에는 되돌리는 치환이 있다. 이러한 연산 규칙이 바로 군의 공리로 추상화된다.

\section{구조의 사슬}

갈루아 이론에 이르는 주요 구조는 다음과 같이 연결된다.

\begin{enumerate}
  \item \textbf{군}: 근들을 서로 바꾸는 대칭을 기록한다.
  \item \textbf{환}: 정수와 다항식처럼 덧셈과 곱셈이 가능한 대상을 다룬다.
  \item \textbf{체}: 0이 아닌 원소로 나눗셈까지 가능한 계수의 세계다.
  \item \textbf{체확장}: 원래 체에 방정식의 근을 첨가하여 더 큰 체를 만든다.
  \item \textbf{갈루아군}: 확장체의 연산을 보존하면서 원래 체를 고정하는 대칭들의 군이다.
\end{enumerate}

따라서 군론, 환론, 체론은 서로 떨어진 단원이 아니다. 하나의 질문, 즉 방정식의 해가 어떤 연산으로 표현될 수 있는가를 이해하기 위한 언어들이다.

\section{첫 번째 예제: $x^2-2$}

$\mathbb{Q}$ 안에는 $x^2-2$의 근이 없다. $\sqrt{2}$를 첨가하면
\[
 \mathbb{Q}(\sqrt{2})=\{a+b\sqrt{2}:a,b\in\mathbb{Q}\}
\]
를 얻는다. 이 체에는 두 근 $\sqrt{2}$와 $-\sqrt{2}$가 모두 들어 있다.

$\mathbb{Q}$의 원소를 그대로 두면서 $\sqrt{2}$를 $-\sqrt{2}$로 보내는 대응
\[
 \sigma(a+b\sqrt{2})=a-b\sqrt{2}
\]
은 덧셈과 곱셈을 보존한다. 가능한 대칭은 항등대응과 $\sigma$ 두 개뿐이며,
\[
 \sigma^2=\mathrm{id}
\]
이다. 이 두 대칭은 원소가 두 개인 순환군을 이룬다.

이 예제는 갈루아 이론의 핵심을 축소해 보여 준다. 방정식의 근을 포함하는 체를 만들고, 원래 계수를 고정하는 그 체의 대칭을 조사한다.

\section{확인 문제}

\begin{exercise}
일반적인 오차방정식에 근호 공식이 없다는 명제와, 특정 오차방정식이 근호로 풀릴 수 있다는 사실이 모순되지 않는 이유를 설명하라.
\end{exercise}

\begin{exercise}
$\mathbb{Q}(\sqrt{3})$에서 $a+b\sqrt{3}$을 $a-b\sqrt{3}$으로 보내는 대응이 덧셈과 곱셈을 보존함을 직접 확인하라.
\end{exercise}

\begin{exercise}
세 개의 기호를 순열하는 모든 방법의 수를 구하고, 두 순열의 합성이 일반적으로 교환법칙을 만족하지 않는 예를 찾아라.
\end{exercise}

\section*{장 요약}

현대대수학은 계산을 포기한 학문이 아니라 계산 가능성의 조건을 연구하는 학문이다. 방정식의 근들 사이에는 대칭이 존재하고, 그 대칭은 군을 이룬다. 근을 담기 위해 체를 확장하고 그 확장체의 대칭을 조사하는 것이 갈루아 이론의 기본 전략이다.
